\section{Frequency-Domain and Finite-Horizon Evidence}
\label{sec:response}

The structural results of Section~\ref{sec:analysis} are visible in two frequency-domain plots, computed before any benchmark. Figure~\ref{fig:transfer-mag} compares $|G(i\omega)|$ on the prey channel at $A=0.25$ for the four mechanisms; Figure~\ref{fig:transfer-phase} shows the phase, which on the collocated channel converges to $-\alpha\pi/2\approx-76.5^{\circ}$ for $\alpha=0.85$ --- the scale-free fractional signature of Theorem~\ref{thm:T4}.

\begin{figure}[htbp]\centering
\includegraphics[width=0.75\textwidth]{fig01_transfer_magnitude}
\caption{$|G(i\omega)|$ on the prey channel at $A=0.25$ for ODE, Caputo ($\alpha=0.85$), DDE ($\tau=0.35$), and latent ($m=3$). The Caputo curve separates over the design band and approaches its $O(\omega^{-\alpha})$ decay.}
\label{fig:transfer-mag}
\end{figure}

\begin{figure}[htbp]\centering
\includegraphics[width=0.90\textwidth]{fig02_phase_2panel}
\caption{Phase on the prey channel. Left: the four mechanisms across frequency. Right: convergence of the Caputo phase to the asymptote $-\alpha\pi/2$, independent of scale.}
\label{fig:transfer-phase}
\end{figure}

\subsection{From exact signatures to finite-horizon ambiguity}
The two panels must be read jointly with Theorem~\ref{thm:T9b}. The $O(\omega^{-\alpha})$ decay and the $-\alpha\pi/2$ phase identify the Caputo channel asymptotically, but an experiment observes a finite band $[\omega_{\min},\omega_{\max}]$, a finite window $[0,T]$, and noise of level $\sigma$. On that restricted set a latent model with enough modes matches the observable part of the response while remaining analytically distinct --- the coexistence of exact separation (Theorem~\ref{thm:T4}) and finite-horizon mimicry (Theorem~\ref{thm:T9b}) stated in frequency-domain terms.

The prey-response computation restates the point in the instrumented channel. Table~\ref{tab:state-approx} lists the certified $L^1$ bounds $\widehat E_m^{\rm state}$ and the induced testing floors at $\alpha=0.85$, $T=12$, $\sigma=0.10$: from $m=4$ to $m=32$ the enclosure shrinks by a factor of $77$ (from $0.5334$ to $0.0069$ at $A=0.25$) and the testing bound rises from $0.374$ to $0.498$, i.e., to within $0.002$ of the equal-prior chance level. Two bounds are reported. Since the two hypotheses are Gaussian with common covariance $\sigma^2I$, the exact minimax error $\Phi(-d/2)$ with $d=\|\Delta\mu\|_2/\sigma$ is available and is uniformly tighter than the Pinsker bound; the difference matters only at small $m$, where Pinsker understates the floor by $0.034$ at $m=4$. We quote the exact bound and retain Pinsker for comparability with the literature.

\begin{table}[htbp]
\centering
\caption{Certified finite-state approximation of the fractional prey response and induced testing lower bounds at $\alpha=0.85$, $T=12$, pulse $\|u\|_2\le0.120$, and $\sigma=0.10$. The exact bound is $\Phi(-d/2)$ with $d=\|\Delta\mu\|_2/\sigma$, valid because the two hypotheses share the covariance $\sigma^2I$; Pinsker is reported for comparison and is uniformly looser.}
\label{tab:state-approx}
\small
\begin{tabular}{ccccccc}
\hline
& \multicolumn{3}{c}{$A=0.25$} & \multicolumn{3}{c}{$A=0.30$}\\
$m$ & $\widehat E_m^{\rm state}$ & $P_e^*$ (exact) & $P_e^*$ (Pinsker) & $\widehat E_m^{\rm state}$ & $P_e^*$ (exact) & $P_e^*$ (Pinsker)\\
\hline
4  & 0.5334 & 0.3745 & 0.340 & 0.4706 & 0.3888 & 0.359\\
8  & 0.1254 & 0.4700 & 0.462 & 0.0898 & 0.4785 & 0.473\\
16 & 0.0273 & 0.4935 & 0.492 & 0.0204 & 0.4951 & 0.494\\
32 & 0.0069 & 0.4983 & 0.498 & 0.0066 & 0.4984 & 0.498\\
\hline
\end{tabular}
\end{table}

Because Table~\ref{tab:state-approx} stops at $m=32$ it shows only part of the obstruction. Table~\ref{tab:state-approx-large} continues it to $m=128$ using an explicit positive-weight exponential-sum surrogate constructed as in Lemma~\ref{thm:T9b} rather than the interval-optimised class $\mathcal G_m$. Because that surrogate is one admissible element and not the infimum, its error exceeds $\widehat E_m^{\rm state}$ --- by factors $2.4$ at $m=4$ and $3.3$ at $m=32$ where the two overlap --- so the induced floors are valid and conservative. They reach $0.4996$ at $m=64$ and $0.49999$ at $m=128$: at attainable latent orders the finite-state rival is indistinguishable from the fractional model under this protocol, and no experiment restricted to the declared budget can separate them.

\begin{table}[htbp]
\centering
\caption{Prey-response $L^1$ error of an explicit positive-weight exponential-sum surrogate (not the certified infimum, hence an upper bound on $\widehat E_m^{\rm state}$) and the induced exact testing floor, $\alpha=0.85$, $A=0.25$, $T=12$, $\|u\|_2\le0.120$, $\sigma=0.10$.}
\label{tab:state-approx-large}
\small
\begin{tabular}{cccc}
\hline
$m$ & $\|g_\alpha-g_m\|_{L^1(0,T)}$ & quadrature error & $P_e^*$ (exact)\\
\hline
4   & $1.29\cdot10^{0}$   & $2.1\cdot10^{-7}$ & 0.2202\\
8   & $4.98\cdot10^{-1}$  & $1.0\cdot10^{-7}$ & 0.3827\\
16  & $1.30\cdot10^{-1}$  & $7.6\cdot10^{-9}$ & 0.4689\\
32  & $2.31\cdot10^{-2}$  & $3.0\cdot10^{-8}$ & 0.4945\\
64  & $1.58\cdot10^{-3}$  & $2.1\cdot10^{-8}$ & 0.4996\\
128 & $4.83\cdot10^{-5}$  & $6.3\cdot10^{-9}$ & 0.49999\\
\hline
\end{tabular}
\end{table}

The table is one-sided. A large lower bound on $P_e^*$ proves that the specified experiment cannot reliably separate $g_\alpha$ from at least one admitted rival $g_m$; a small bound proves nothing in the converse direction, so such cells are inconclusive rather than certified-easy. Figure~\ref{fig:atlas} applies this reading across $(\alpha,A,m)$: cells outside the focal Matignon-stable regime ($\alpha>\alpha^*(A)$) are hatched and excluded from the headline interpretation.

\begin{figure}[htbp]\centering
\includegraphics[width=0.92\textwidth]{fig18_safe_discrimination_atlas}
\caption{Protocol-restricted discrimination atlas over $(\alpha,A,m)$: lower bound on the model-selection error under the declared pulse protocol. The bound increases with latent budget $m$ and toward weakly separated regimes; hatched cells violate the Matignon condition $\alpha<\alpha^*(A)$ and are excluded from stable-regime conclusions.}
\label{fig:atlas}
\end{figure}
